\documentclass[12pt]{beamer}
\usetheme{Malmoe}
\usecolortheme{dolphin}
%\usefonttheme{professionalfonts}
\useoutertheme{miniframes} 
%{infolines}
%\usepackage{xmpmulti}
%\usepackage{fouriernc}
%\usepackage[T1]{fontenc}
%\usepackage{fourier}
%\usepackage[scaled=0.9]{helvet}

\usepackage{bm}
\usepackage{titlesec}
\usepackage{graphicx,subfigure}
\usepackage{amsmath}
\usepackage{amsfonts}

%\usepackage{mathptmx}
%\usepackage[scaled=.90]{helvet}
%\usepackage{courier}
%\usepackage[T1]{fontenc}


\usepackage{mathpazo}

%\usepackage[bitstream-charter]{mathdesign}
\usepackage{amssymb}
\usepackage{url}
\usepackage{appendix}
\usepackage{indentfirst}
\usepackage{makeidx}
\usepackage{array}


\newcommand{\proposition}[1]{\bigskip\noindent\fbox{\begin{minipage}{0.95\textwidth}~\begin{minipage}{0.95\textwidth}#1\end{minipage}\\\smallskip\end{minipage}}\bigskip\vline\vline\vline}

\newcommand{\excercise}{\addtocounter{subsection}{+1}
\subsubsection{Ex~\underline{\thesubsection}}}
\newcommand{\stag}[1]{\tag{\arabic{equation}#1}}

\newcommand{\qdo}[0]{\quad}
\newcommand{\qdt}[0]{\quad\quad}
\newcommand{\qdthr}[0]{\quad\quad\quad}
\newcommand{\qdf}[0]{\quad\quad\quad\quad}
\newcommand{\qdfiv}[0]{\quad\quad\quad\quad\quad}

\newcommand{\norm}[1]{\|#1\|}
\newcommand{\sdelta}{{\scriptstyle\Delta}}
\newcommand{\bra}[1]{\langle#1|}
\newcommand{\ket}[1]{|#1\rangle }
\newcommand{\braket}[2]{\langle#1|#2\rangle }
\newcommand{\proj}[2]{|#1\rangle\langle#2|}
\newcommand{\expect}[2]{\langle #1\rangle_{#2}}
\newcommand{\Kappa}[0]{\bm{\mathcal{K}}}
\newcommand{\parity}[0]{\bm{\mathcal{P}}}
\newcommand{\parityl}[0]{\bm{p}}
\newcommand{\charge}[0]{\bm{\mathcal{C}}}
\newcommand{\ptplus}[0]{\\[+5pt]}
%%%%%%

\title{Lorentz Symmetry, Weyl Spinors, Chirality and Dirac Equation}
\author{Kow Lung Chang}
\institute{Physics Department, National Taiwan University}
%%%%%%

\begin{document}
\begin{frame}
\titlepage
\end{frame}

%%%%%page1
\begin{frame}
\frametitle{Contents}
$\blacklozenge$ Minkowski Space and Lorentz Transformation \ptplus

$\blacklozenge$ Generators of Lorentz Group\ptplus

$\blacklozenge$ Irreducible Representations of Lorentz Group \\ ~~~and Weyl 
Spinors \ptplus

$\blacklozenge$ SO(3,1) and SL(2,C)\ptplus

$\blacklozenge$ Chiral Transformation and Spinor Algebra\ptplus

$\blacklozenge$ Spinor space and Co-spinor space\ptplus

$\blacklozenge$ Dirac spinor and Dirac equation\ptplus

$\blacklozenge$ Invariance of the $ \gamma $ matrices in all Lorentz frames\ptplus

$\blacklozenge$ Zero Mass Limit and Helicity of Weyl spinors\ptplus

\end{frame}

\begin{frame}
\frametitle{Minkowski Space and Lorentz Transformation}

Difine the Minkowski contravariant 4-vector as:

\begin{align}
x^\mu=(x^0=ct, \vec{x}),
\end{align}

\noindent and the Minkowski covariant 4-vector as:

\begin{align}
x_\mu=(x^0=-ct, \vec{x}),
\end{align}

\noindent with the metric tensor

\begin{align}
g_{\mu\nu}=0~~~\mbox{if}~~~ \mu\neq\nu~;~~~-g_{00}=g_{11}=g_{22}=g_{33}=1.
\end{align}
\end{frame}

\begin{frame}
A linear transformation on the $ x^\mu $ given as follows

\begin{align}
x'^{\mu}=\bm{\Lambda}^\mu_\nux^\nu
\end{align}

\noindent is called Homogeneous Lorentz transformation ( HLT ), or simply LT if the following condition is met: 

\begin{align}
x'^{\mu}x'_\mu=x^\nux_\mu
\end{align}

\noindent or in matrix notation as

\begin{align}\label{transform-41}
\bm{\Lambda}^T\bm{g}\bm{\Lambda}=\bm{g} ~~~\mbox{ or }~~~ \bm{\Lambda}^T\bm{g}=\bm{g}\bm{\Lambda}^{-1} .
\end{align}

\end{frame}

\begin{frame}
Since there exists an identity Lorentz transformation,$ \Lambda=\mathbf{I} $, and an inverse Lorentz transformation, $ \bm{\Lambda}^{-1} $, namely both $ \mathbf{I} $ and $ \bm{\Lambda}^{-1} $ exist.

Therefore LT forms a group SO(3,1) because:  

\begin{align}\label{transform-41}
(\bm{\Lambda}_1\bm{\Lambda}_2)^Tg(\bm{\Lambda}_1\bm{\Lambda}_2)= \bm{\Lambda}_2^T\bm{\Lambda}_1^Tg\bm{\Lambda}_1\bm{\Lambda}_2=g.
\end{align}

The condition $ \mbox{det} \bm{\Lambda}=1 $ is automatically satisfied. We shall only 
consider the proper LT in which $ \bm{\Lambda}^0_0\geqslant 1 $ in this lecture.
\end{frame}

\begin{frame}
Since the condition $ \bm{\Lambda}^tg\bm{\Lambda}=g $ provide 10 constraints among 16 matrix elements of $ \bm{\lambda} $, the remaining 6 independent coefficients serve as the 6 group parameters, specified as $ \bm{\Lambda}=\bm{\Lambda}(\vec{\theta},\vec{\xi}) $ and

\addtocounter{equation}{+1}
\begin{align*}
\vec{\theta}&=(\theta^1,\theta^2,\theta^3)=\mbox{ rotation},\stag{a}\\
\vec{\xi}&=(\xi^1,\xi^2,\xi^3)=\mbox{ Lorentz boost}\stag{b}.
\end{align*}
\end{frame}

\begin{frame}
\frametitle{Generators of Lorentz Group}

The generators of the group are given:

\begin{align}
A_i=\frac{\partial}{\partial\theta^i}\bm{\Lambda}(\vec{\theta},\vec{\xi}),~~~B_i=\frac{\partial}{\partial\xi^i}\bm{\Lambda}(\vec{\theta},\vec{\xi}),\stag{a,b}
\end{align}

For the Lorentz boost along 1-axis with angle $ \xi $,

\begin{align}
\bm{\Lambda}=\begin{pmatrix}
\cosh\xi&-\sinh\xi&0&0\\
-\sinh\xi&\cosh\xi&0&0\\
0&0&1&0\\
0&0&0&1\end{pmatrix}.
\end{align}

\noindent where $ \xi=\tanh^{-1}\beta $ and
\end{frame}

\begin{frame}
\addtocounter{equation}{1}
\begin{align*}
B_1=\begin{pmatrix}
0&-1&0&0\\
-1&0&0&0\\
0&0&0&0\\
0&0&0&0\end{pmatrix}.~~~~~~\stag{a}
\end{align*}

\noindent similarly,

\begin{align*}
B_2=\begin{pmatrix}
0&0&-1&0\\
0&0&0&0\\
-1&0&0&0\\
0&0&0&0\end{pmatrix},~~~
B_3=\begin{pmatrix}
0&0&0&-1\\
0&0&0&0\\
0&0&0&0\\
-1&0&0&0\end{pmatrix}.\stag{b,c}
\end{align*}
\end{frame}

\begin{frame}
And for the generators if rotation, we have

\addtocounter{equation}{1}
\begin{align*}
A_1=\begin{pmatrix}
0&0&0&0\\
0&0&0&0\\
0&0&0&1\\
0&0&-1&0\end{pmatrix}~~~~~~\stag{a}
\end{align*} 

\noindent and

\begin{align*}
\mbox{and}~~A_2=\begin{pmatrix}
0&0&0&0\\
0&0&0&-1\\
0&0&0&0\\
0&1&0&0\end{pmatrix},~~~
A_3=\begin{pmatrix}
0&0&0&0\\
0&0&1&0\\
0&-1&0&0\\
0&0&0&0\end{pmatrix}.\stag{b,c}
\end{align*}
\end{frame}
\begin{frame}
SO(3,1) Lie algebra as:

\begin{align}
[A_i,A_j]=-\epsilon_{ij}^kA_k,~~~[A_i,B_j]=-\epsilon_{ij}^kB_k,
\end{align}

Canonical formulation of algebra:

\addtocounter{equation}{+1}
\begin{align*}
M^{\mu\nu}=x^{\mu}\frac{\partial}{\partial x_\nu}-x^{\nu}\frac{\partial}{\partial x_\mu},\stag{a}
\end{align*}
\begin{align*}
[M^{\mu\nu},M^{\alpha\beta}]=-g^{\nu\beta}M^{\mu\alpha}-g^{\mu\alpha}M^{\nu\beta}
+g^{\nu\alpha}M^{\mu\beta}+g^{\mu\beta}M^{\nu\alpha}.\stag{b}
\end{align*}
\end{frame}


\begin{frame}
If we denote

\addtocounter{equation}{1}
\begin{align*}
L_i=\frac{1}{2}\left(\frac{A_i}{i}+B_i\right),~~~\mbox{and}~~~R_i=\frac{1}{2}\left(\frac{A_i}{i}-B_i\right).\stag{a,b}
\end{align*}

The algebra takes as

\addtocounter{equation}{1}
\begin{align*}
\notag[L_i,L_j]&=i\epsilon^k_{ij}L_k,\stag{a}\\
\notag[L_i,R_i]&=0,\stag{b}\\
[R_i,R_j]&=i\epsilon^k_{ij}R_k.\stag{c}
\end{align*} 
\end{frame}

\begin{frame}
\frametitle{Irreducible Representations of \\
~~~~~~~~~~~Lorentz Group and Weyl Spinors}

Consider the finite dimensional representations, denoted by $ (l, r) $ with the basis\\[-10pt]

\begin{align}
\ket{l,m}\otimes\ket{r,n}\equiv\ket{l,m;r,n}
\end{align}

\noindent where\\[-10pt]

\begin{align}
-l\leqslant m\leqslant l,~~~-r\leqslant n\leqslant r,~~~\mbox{and}~~~l,r=~\mbox{half integers.}
\end{align} 

The simpliest representation of the generators, an one dimensional
$ (0,0) $-representation read as\\[-10pt]

\begin{align}
\bra{0,0;0,0}L_i \ket{0,0;0,0}=\bra{0,0;0,0}R_i \ket{0,0;0,0}=0,
\end{align}
\end{frame}
\begin{frame}
$ (\frac{1}{2}, 0) $-representation: left-handed-spinor

\begin{align}
\ket{\frac{1}{2},m;0,0}
\end{align}

$ (0, \frac{1}{2}) $-representation: right-handed-spinor

\begin{align}
\ket{0,0;\frac{1}{2},n}
\end{align}

\noindent then we have

\addtocounter{equation}{1}
\begin{align*}
L^{(\frac{1}{2},0)}_i&=\frac{1}{2}\bm{\sigma_i},&R^{(\frac{1}{2},0)}_i&=0,\tag{\arabic{chapter}.\arabic{equation}a}\\
L^{(0,\frac{1}{2})}_i&=0,&R^{(0,\frac{1}{2})}_i&=\frac{1}{2}\bm{\sigma_i},\tag{\arabic{chapter}.\arabic{equation}b}
\end{align*}

\end{frame}
\begin{frame}
\noindent which lead to

\addtocounter{equation}{1}
\begin{align*}
A^{(\frac{1}{2},0)}_i&=\frac{i}{2}\bm{\sigma_i},&B^{(\frac{1}{2},0)}_i&=\frac{1}{2}\bm{\sigma_i},\tag{\arabic{chapter}.\arabic{equation}a}\\
A^{(0,\frac{1}{2})}_i&=\frac{i}{2}\bm{\sigma_i},&B^{(0,\frac{1}{2})}_i&=-\frac{1}{2}\bm{\sigma_i},\tag{\arabic{chapter}.\arabic{equation}b}
\end{align*}

\noindent and the 2-dimensional irreducible representation of Lorentz group as

\addtocounter{equation}{1}
\begin{align*}
D^{(\frac{1}{2},0)}(\vec{\theta},\vec{\xi})=\exp\left(\frac{i}{2}\vec{\sigma}\cdot(\vec{\theta}-i\vec{\xi})\right),\stag{a}\label{dspin1}\\
D^{(0,\frac{1}{2})}(\vec{\theta},\vec{\xi})=\exp\left(\frac{i}{2}\vec{\sigma}\cdot(\vec{\theta}+i\vec{\xi})\right).\stag{b}
\end{align*}
\end{frame}
\begin{frame}

As a quick check that

\addtocounter{equation}{+1}
\begin{align*}
D^{(\frac{1}{2},0)\dag}(\vec{\theta},\vec{\xi})=\exp\left(-\frac{i}{2}\vec{\sigma}\cdot(\vec{\theta}+i\vec{\xi})\right)\neq D^{(\frac{1}{2},0)}(\vec{\theta},\vec{\xi})^{-1},
\end{align*}

\begin{align*}
D^{(0,\frac{1}{2})\dag}(\vec{\theta},\vec{\xi})=\exp\left(-\frac{i}{2}\vec{\sigma}\cdot(\vec{\theta}-i\vec{\xi})\right)\neq D^{(0,\frac{1}{2})}(\vec{\theta},\vec{\xi})^{-1},
\end{align*}

Let us perform the identifications

\addtocounter{equation}{1}
\begin{align*}
\ket{\frac{1}{2},\frac{1}{2}}_l\longmapsto\begin{pmatrix}1\\0\end{pmatrix}=e_1,~~~
\ket{\frac{1}{2},-\frac{1}{2}}_l\longmapsto\begin{pmatrix}0\\1\end{pmatrix}=e_2,\stag{a,b}
\end{align*}
\end{frame}


\begin{frame}
\noindent then\\[-10pt]

\begin{align}
\psi_l(x)=\psi^a_l(x)e_a=\begin{pmatrix}\psi^1_l(x)\\\psi^2_l(x)\end{pmatrix},
\end{align}

the Lorentz Transformation as

\begin{align}
\psi_l(x)\longmapsto\psi_l'(x')=D^{(\frac{1}{2},0)}(\vec{\theta},\vec{\xi})\psi_l(\bm{\Lambda}^{-1}x'),
\end{align}

Similarly if\\[-10pt] 

\begin{align}
\psi_r(x)=\psi^a_r(x)f_a=\psi^1_r(x)\begin{pmatrix}1\\0\end{pmatrix}+ \psi^2_r(x)\begin{pmatrix}0\\1\end{pmatrix}=\begin{pmatrix}\psi^1_r(x)\\\psi^2_r(x)\end{pmatrix},
\end{align}

\noindent the transformation reads as\\[-10pt]

\begin{align}
\psi_r(x)\longmapsto\psi_r'(x')=D^{(0,\frac{1}{2})}(\vec{\theta},\vec{\xi})\psi_r(\bm{\Lambda}^{-1}x')).
\end{align}
\end{frame}
\begin{frame}
\frametitle{SO(3,1) and SL(2,C)}

SL(2,C) transformation in $ \mathcal{C}^2 $-space:

\begin{align}
\xi'=\begin{pmatrix}\xi'^1\\\xi'^2\end{pmatrix}=\bm{L}\begin{pmatrix}\xi^1\\\xi^2\end{pmatrix}= \begin{pmatrix}a&b\\c&d\end{pmatrix}\begin{pmatrix}\xi^1\\\xi^2\end{pmatrix}.
\end{align}

\noindent where

\begin{align}
\mbox{det}~\bm{L}=\left|\begin{array}{cc}a&b\\c&d\end{array}\right|=ab-bc=1,
\end{align}

If we exponentiate $ L $ by a $ 2\times 2 $ matrix $ \bm{A} $, i.e.

\begin{align}
\bm{L}=e^{\bm{A}},
\end{align}
\end{frame}

\begin{frame}
Then we have the following proposition

\proposition{
\subsection*{\textit{Proposition 1.}}
{\it If a matrix} $ \bm{L} $ {\it can be expressed as} $ \bm{L}=e^{\bm{A}} $, {\it then }

\begin{align}
\mbox{det}~\bm{L}=e^{\mbox{Tr}~\bm{A}}. 
\end{align}
}

Hence we ensure that

\begin{align}
\mbox{det}~D^{(\frac{1}{2},0)}(\vec{\theta},\vec{\xi})= \mbox{det}~e^{\frac{i}{2}\vec{\sigma}\cdot(\vec{\theta}-i\vec{\xi})}=e^{\mbox{Tr}~\frac{i}{2}\vec{\sigma}\cdot(\vec{\theta}-i\vec{\xi})}=1,
\end{align}

\noindent and

\begin{align}
\mbox{det}~D^{(0,\frac{1}{2})}(\vec{\theta},\vec{\xi})= \mbox{det}~e^{\frac{i}{2}\vec{\sigma}\cdot(\vec{\theta}+i\vec{\xi})}=e^{\mbox{Tr}~\frac{i}{2}\vec{\sigma}\cdot(\vec{\theta}+i\vec{\xi})}=1.
\end{align}

\end{frame}

\begin{frame}
The isomorphism of SL(2,C) onto SO(3,1) in Lorentz transformation can be demonstrated as follows:

let

\begin{align}
\bm{X}=x^\mu\bm{\sigma}_\mu=\begin{pmatrix}-x^0+x^3&x^1-ix^2\\x^1+ix^2&-x^0-x^3\end{pmatrix},
\end{align}

\noindent  and

\begin{align}
\mbox{det}~\bm{X}=(x^0)^2-\vec{x}^2.
\end{align}

\noindent which leads to the Lorentz Transformation on $ \bm{X} $ as

\begin{align}
\bm{X}'=D^{(\frac{1}{2},0)}(\vec{\theta},\vec{\xi})\bm{X}D^{(\frac{1}{2},0)\dag}(\vec{\theta},\vec{\xi}),
\end{align}

\end{frame}

\begin{frame}
\noindent because of the invariance of the length of the space-time vector,\\[-10pt]

\begin{align}
\mbox{det}~\bm{X}'=\mbox{det}~\bm{X}.
\end{align}

As an example when $ \mathcal{O}' $-frame is boost along the 3rd axis, i.e.

\begin{align}
\bm{X}'=e^{\frac{1}{2}\bm{\sigma}_3\xi}\bm{X}e^{\frac{1}{2}\bm{\sigma}_3\xi}=
\begin{pmatrix}e^{\frac{1}{2}\xi}&0\\0&e^{-\frac{1}{2}\xi}\end{pmatrix}\bm{X}
\begin{pmatrix}e^{\frac{1}{2}\xi}&0\\0&e^{-\frac{1}{2}\xi}\end{pmatrix}.
\end{align}

\noindent we regain the LT as follows

\addtocounter{equation}{1}
\begin{align*}
x'^0&=\cosh\xi x^0-\sinh\xi x^3,\stag{a}\\
x'^1&=x^1,\stag{b}\\
x'^2&=x^2,\stag{c}\\
x'^3&=-\sinh\xi x^0+\cosh\xi x^3.\stag{d}
\end{align*}

\end{frame}
\begin{frame}
\frametitle{Chiral Transformation}
$ \Kappa =$ chiral operator, which is a discrete transformation between left handed irreducible representations and the right handed irreducible representations

It is an antilinear operator, i.e.
 
\begin{align}
\Kappa(a\psi+b\varphi)=a^*\Kappa\psi+b^*\Kappa\varphi,
\end{align}

\noindent as well as an antiunitary operator:

\begin{align}
(\Kappa\psi,\Kappa\varphi)=(\varphi,\psi)=(\psi,\varphi)^*.
\end{align} 
\end{frame}

\begin{frame}
It is nothing to with the space-time coordinates, hence 


\begin{align}
\Kappa A_i\Kappa^{-1}=A_i,~~~\Kappa B_i\Kappa^{-1}=B_i.
\end{align}

\noindnet but the operators $ L_i $ and $ R_i $ transform as follows

\addtocounter{equation}{1}
\begin{align*}
\Kappa L_i\Kappa^{-1}=\frac{1}{2}\Kappa\left(\frac{A_i}{i}+B_i\right)\Kappa^{-1}=-R_i,~~~
\Kappa R_i\Kappa^{-1}=-L_i,\stag{a,b}
\end{align*}

\noindent therefore we reach the following proposition if the basis of $ (\frac{1}{2},0) $- and $ (0,\frac{1}{2}) $-representation are abbreviated by

\addtocounter{equation}{1}
\begin{align*}
\ket{j,m;0,0}=L_{jm},\stag{a}\\
\ket{0,0,k,n}=R_{kn},\stag{b}
\end{align*}
\end{frame}

\begin{frame}
\noindent then

\proposition{
\subsection*{\textit{Proposition 2.}}

\it The vector $ \Kappa L_{jm} $ is the eigenvector of $ R^2 $ and $ R_3 $     with the eigenvalues $ j(j+1) $ and $ -m $ respectively. While the vector $ \Kappa R_{kn} $ is the eigenvector of $ L^2 $ and $ L_3 $ with the eigenvalues $ k(k+1) $ and $ -n $ respectively.
}

Since $ \KappaL^2=R^2\Kappa $, $ \Kappa\vec{L}=-\vec{R}\Kappa $

\end{frame}
\begin{frame}
\noindent then we have \\[-12pt]

\begin{align*}
R^2\Kappa L_{jm}=\Kappa L^2L_{jm}=j(j+1)\Kappa L_{jm},\\
R_3\Kappa L_{jm}=-\Kappa L_3L_{jm}=-m\Kappa L_{jm},
\end{align*}

Therefore \\[-12pt]

\begin{align*}
\Kappa L_{jm}=\gamma(m)R_{j -m}.
\end{align*}

Similarly \\[-18pt]

\begin{align*}
\Kappa R_{kn}=\delta{n}L_{k -n}.
\end{align*}

Hence we have \\[-18pt]

\begin{align} 
(\Kappa L_{jm},\Kappa L_{jm'})=\gamma^*(m)\gamma(m')(R_{j,-m}, R_{j,-m'})=(L_{jm'},L_{jm}),
\end{align}

\noindent or $ \gamma^*(m)\gamma(m')\delta_{-m,-m'}=\delta_{m'm} $

\end{frame}
\begin{frame}
Take spinor for instance, we obtain that

\addtocounter{equation}{1}
\begin{align*}
\Kappa L_{\frac{1}{2} m}=\gamma(n)\Kappa R_{\frac{1}{2} -m}, ~~~~~
\Kappa R_{\frac{1}{2} n}=\delta(n)\Kappa L_{\frac{1}{2} -n}, \stag{ab}
\end{align*}

\noindent or

\addtocounter{equation}{1}
\begin{align*}
\Kappa \begin{pmatrix}1\\0\end{pmatrix}_l=\gamma(\frac{1}{2})\begin{pmatrix}0\\-1\end{pmatrix}_{r}~~~~~~~~
\Kappa \begin{pmatrix}0\\1\end{pmatrix}_l=\gamma(-\frac{1}{2})\begin{pmatrix}-1\\0\end{pmatrix}_{r},  \stag{a,b}
\end{align*}

\begin{align*}
\Kappa \begin{pmatrix}1\\0\end{pmatrix}_r=\delta(\frac{1}{2})\begin{pmatrix}0\\-1\end{pmatrix}_{l}~~~~~~~~
\Kappa \begin{pmatrix}0\\1\end{pmatrix}_r=\delta(-\frac{1}{2})\begin{pmatrix}-1\\0\end{pmatrix}_{l},  \stag{c,d}
\end{align*}
\end{frame}

\begin{frame}
Let us evaluate the following matrix elements

\begin{align}
\notag D^{(\frac{1}{2},0)}_{mm'}&=(L_{\frac{1}{2},m},e^{\vec{\theta}\cdot\vec{A}+\vec{\xi}\cdot\vec{B}}L_{\frac{1}{2},m'})
=(\Kappa e^{\vec{\theta}\cdot\vec{A}+\vec{\xi}\cdot\vec{B}}L_{\frac{1}{2},m'},\Kappa L_{\frac{1}{2},m})\\
\notag &=\gamma^*(m')\gamma(m)(e^{\vec{\theta}\cdot\vec{A}+\vec{\xi}\cdot\vec{B}}R_{\frac{1}{2},-m'},R_{\frac{1}{2},-m})\\&=\gamma(m)D^{(0,\frac{1}{2})*}_{-m,-m'}\gamma^*(m'),
\end{align}

\noindent or

\begin{align}
D^{(\frac{1}{2},0)}=\begin{pmatrix}0&\gamma(\frac{1}{2})\\\gamma(-\frac{1}{2})&0\end{pmatrix}D^{(0,\frac{1}{2})*}\begin{pmatrix}0&\gamma^*(-\frac{1}{2})\\\gamma^*(\frac{1}{2})&0\end{pmatrix}.
\end{align}
\end{frame}

\begin{frame}
Similarly,

\begin{align}
D^{(0,\frac{1}{2})}=\begin{pmatrix}0&\delta(\frac{1}{2})\\\delta(-\frac{1}{2})&0\end{pmatrix}D^{(\frac{1}{2},0)*}\begin{pmatrix}0&\delta^*(-\frac{1}{2})\\\delta^*(\frac{1}{2})&0\end{pmatrix}.
\end{align}

In order to be consistent with the LT for spinors, one chooses

\begin{align}
\gamma(\frac{1}{2})=\delta(\frac{1}{2})=-\gamma(-\frac{1}{2})=-\delta(-\frac{1}{2})=1
\end{align}

\noindent hence we have\\[-10pt]

\begin{align}
\epsilon=\begin{pmatrix}0&1\\-1&0\end{pmatrix}
\end{align}

\noindent and\\[-10pt]

\begin{align}
\epsilon\bm{\sigma}_i^*\epsilon^{-1}=\begin{pmatrix}0&1\\-1&0\end{pmatrix} \bm{\sigma}_i^* \begin{pmatrix}0&-1\\1&0\end{pmatrix}=-\bm{\sigma}_i.
\end{align}
\end{frame}


\begin{frame}
\frametitle{Spinor space and Co-spinor space}

\begin{align}
\mbox{Let }\left\{\begin{array}{ll}\mathcal{V}_{(\frac{1}{2},0)}&:\mbox{ left-hand spinor space }\\
e_{a}&: (a=1, 2) \mbox{ two left-handed spinor }\end{array}\right.
\end{align}

A spinor in $ \mathcal{V}_{(\dot{\frac{1}{2}},0)} $

\begin{align}
\left\{\begin{array}{ll}\mathcal{V}_{(\dot{\frac{1}{2}},0)}&:\mbox{ co-left-hand spinor space }\\
e_{\dot{a}}&: (a=1, 2) \mbox{ two co-left-handed spinor }\end{array}\right.
\end{align}

\noindent which is related to the left-handed spinor by

\begin{align}
\dot{\psi}=\psi^T\epsilon^T=(\dot{\psi}_1,\dot{\psi}_1)=(\psi^1,\psi^2)\epsilon^T=(\psi^1,\psi^2)\epsilon^{-1}.
\end{align}

\end{frame}

\begin{frame}
\noindent and the corresponding LT is given as

\begin{align}
\dot{\psi}\stackrel{L.T.}{\longmapsto}\dot{\psi}'&=\psi'^T\epsilon^{-1}=
\psi^TD^{T(\frac{1}{2},0)}\epsilon^{-1}\\
&=\dot{\psi}(\epsilon D^{(\frac{1}{2},0)*}\epsilon^{-1})^{\dag}= \dot{\psi}D^{(0,\frac{1}{2})\dag}.
\end{align}

Similarly the LT for the co-right-handed spinor is given as

\begin{align}
\dot{\varphi}\stackrel{L.T.}{\longmapsto}\dot{\varphi}'=\varphi'^T\epsilon^T=
\dot{\varphi}D^{(\frac{1}{2},0)\dag}.
\end{align}

Let us construct the 4-dimensional product space as

\addtocounter{equation}{+1}
\begin{align*}
\mathcal{V}_{(\frac{1}{2},\dot{\frac{1}{2}})}:~e_a\dot{f}^b~~\mbox{as basis}\stag{a}\\
\mathcal{V}_{(\dot{\frac{1}{2}},\frac{1}{2})}:~\dot{e}_af^b~~\mbox{as basis}\stag{b}
\end{align*}

\end{frame}

\begin{frame}
then any element in $ \mathcal{V}_{(\frac{1}{2},\dot{\frac{1}{2}})} $ and in $ \mathcal{V}_{(\dot{\frac{1}{2}},\frac{1}{2})} $ can be written respectively as\\[-15pt]

\begin{align}
U^{(\frac{1}{2},\dot{\frac{1}{2}})} =\begin{pmatrix}u^1_{\dot{1}}&u^1_{\dot{2}}\\u^2_{\dot{1}}&u^2_{\dot{2}}\end{pmatrix}.
\end{align}

\noindent and\\[-15pt]

\begin{align}
U^{(\dot{\frac{1}{2}},\frac{1}{2})} =\begin{pmatrix}u^{\dot{1}}_1&u^{\dot{2}}_1\\u^{\dot{1}}_2&u^{\dot{2}}_2\end{pmatrix}.
\end{align}

It is obvious that the space-time matrix

\begin{align}
\bm{X}=x^\mu\bm{\sigma}_\mu=\begin{pmatrix}-x^0+x^3&x^1-ix^2\\x^1+ix^2&-x^0-x^3\end{pmatrix},
\end{align}

\noindent transforms as an element of $ \mathcal{V}_{(\frac{1}{2},\dot{\frac{1}{2}})} $-representation, while

\begin{align}
\bm{X}_c=\epsilon\bm{X}^*\epsilon^{-1}
\end{align}

\noindent transforms as an element of $ \mathcal{V}_{(\dot{\frac{1}{2}},\frac{1}{2})} $-representation.

\end{frame}

\begin{frame}
We are now in the position to emphasize the next proposition

\proposition{
\subsection*{\textit{Proposition 3.}}
{\it An operator of the} $ (\frac{1}{2},\dot{\frac{1}{2}}) $-{\it representation acts upon a vector of the} $ (0,\frac{1}{2}) $-{\it representation yields a vector of} $ (\frac{1}{2},0) $-{\it representation. Conversely an operator of the} $ (\dot{\frac{1}{2}},\frac{1}{2}) $-{\it representation acts upon a vector of the} $ (\frac{1}{2},0) $-{\it representation will yield a vector of the} $ (0,\frac{1}{2}) $-{\it representation.}
}

The proof goes as; 
If $ \bm{A}\in (\frac{1}{2},\dot{\frac{1}{2}}) $-representation, 
$ \xi\in(0,\frac{1}{2}) $-representation, then \\[-15pt]

\begin{align}
\eta\stackrel{L.T.}{\longmapsto}\eta'=\bm{A}'\xi'= D^{(\frac{1}{2},0)}\bm{A}D^{(\frac{1}{2},0)\dag}D^{(0,\frac{1}{2})}\xi=D^{(\frac{1}{2},0)}\eta,
\end{align}
\end{frame}


\begin{frame}
Hence we have a left-handed spinor, i.e.

\begin{align}
\eta=\bm{A}\xi,~~~~~~~~~~~~~
\end{align}

Similarly that

\begin{align*}
\mbox{if }\left\{\begin{array}{ll}\bm{A}&\in (\dot{\frac{1}{2}},\frac{1}{2})-\mbox{representation},\\[+8pt]
\eta&\in(\frac{1}{2},0)-\mbox{representation}.\end{array}\right.
\end{align*}

Therefore we reach as follows,

\begin{align}
\xi\stackrel{L.T.}{\longmapsto}\xi'=\bm{A}_c'\eta'= D^{(0,\frac{1}{2})}\bm{A}_cD^{(0,\frac{1}{2})\dag}D^{(\frac{1}{2},0)}\xi=D^{(0,\frac{1}{2})}\xi,
\end{align}
\end{frame}


\begin{frame}
\frametitle{Dirac spinor and Dirac equation}
Since any Lorentz 4-vector with the construction      

\begin{align}
U=\sigma_\mu u^\mu,~~~~ U_c=\sigma^c_\mu u^\mu,
\end{align}

transforms as $ (\frac{1}{2},\dot{\frac{1}{2}}) $-representation and $ (\dot{\frac{1}{2}},\frac{1}{2}) $-representation respectively. Therefore,

\addtocounter{equation}{1}
\begin{align*}
\bm{P}_c\psi_l&=m_0c\psi_r,\tag{\arabic{chapter}.\arabic{equation}a}~~~~~~~~\\
\bm{P}\psi_r&=m_0c\psi_l,\tag{\arabic{chapter}.\arabic{equation}b}~~~~~~~~
\end{align*}

\noindent hence we have

\begin{align}
\begin{pmatrix}0&\bm{P}_c\\\bm{P}&0\end{pmatrix}
\begin{pmatrix}\psi_r\\\psi_l\end{pmatrix}=
m_0c\begin{pmatrix}\psi_r\\\psi_l\end{pmatrix},
\end{align}
\end{frame}

\begin{frame}
 If we define the Dirac spinor $ \psi_d $ as $ \psi_r\oplus\psi_l $, then
 
\begin{align}\label{dirac-eq2}
\begin{pmatrix}0&\bm{P}_c\\\bm{P}&0\end{pmatrix}\psi_d(x)=m_0c\psi_d(x),~~~~~
\end{align}

\noindent or 

\begin{align}
\left[\begin{pmatrix}0&i\bm{\sigma}_c^\mu\\i\bm{\sigma}^\mu&0\end{pmatrix}\partial_\mu+\frac{m_0c}{\hbar}\right]\psi_d(x)=0,
\end{align}

\noindent which can be cast into

\begin{align}\label{dirac-eq1}
\left(i\gamma^\mu\partial_\mu+\frac{m_0c}{\hbar}\right)\psi_d(x)=0,
\end{align}
\end{frame}

\begin{frame}
\noindent with $ \gamma^\mu $ defined as

\begin{align}\label{gamma-mu}
\gamma^\mu=\begin{pmatrix}0&\bm{\sigma}_c^\mu\\\bm{\sigma}^\mu&0\end{pmatrix}~~~
\mbox{or}~~~\gamma^0=\begin{pmatrix}0&\mathbf{I}\\\mathbf{I}&0\end{pmatrix},~~~
\gamma^i=\begin{pmatrix}0&-\bm{\sigma}^i\\\bm{\sigma}^i&0\end{pmatrix}.
\end{align}   

The covariant formulation of Dirac equation does not imply that

\begin{align*}
\gamma_\mu\partial^\mu=\gamma'_\mu\partial'^\mu
\end{align*}

In fact that we have the following proposition.
\end{frame}

\begin{frame}
\proposition{
\subsection*{\textit{Proposition 4.}}
\it The gamma matrices $ \gamma^\mu $ are universal in all Lorentz frame, namely the Dirac equation in another Lorentz frame, i.e. the $ \mathcal{O}' $-system always takes the same gamma matrices $ \gamma^\mu $ used in   $ \mathcal{O} $-system. The equation in $ \mathcal{O}' $-system is expressed as

\begin{align*}
\left(i\gamma^\mu\partial'_\mu+\frac{m_0c}{\hbar}\right)\psi_d'(x')=0.
\end{align*}
}
\end{frame}

\begin{frame}
It can be proved by defining

$$  
D(\vec{\theta},\vec{\xi})= D^{(0,\frac{1}{2})}(\vec{\theta},\vec{\xi}) \oplus D^{(\frac{1}{2},0)}(\vec{\theta},\vec{\xi}) 
$$

\noindent and multiplying it upon the Dirac equation as follows

\begin{align}
\left(iD\gamma^\mu D^{-1}\partial_\mu+\frac{m_0c}{\hbar}\right)D\psi_d(x)=0.
\end{align}

One can show that\\[-13pt] 

\begin{align}
D\gamma^\mu D^{-1}=\bm{\Lambda}^\mu_\nu\gamma^\nu,
\end{align}

\noindent and the Dirac equation in the new Lorentz frame, i.e.   $ \mathcal{O}' $–frame reads as\\[-13pt]

\begin{align}
\left(i\gamma^\mu\partial'_\mu+\frac{m_0c}{\hbar}\right)\psi_d'(
x')=0,
\end{align}

\noindent by identifying $ \psi'_d(x')\equiv D(\vec{\theta},\vec{\xi})\psi_d(x)=D(\vec{\theta},\vec{\xi})\psi_d(\bm{\Lambda}^{-1}
x') $.
\end{frame}

\begin{frame}
\frametitle{Zero Mass Limit and Helicity of Weyl spinors}

Last demonstration for zero mass limit in Dirac equation

\begin{align}
\begin{pmatrix}0&\bm{P}_c\\\bm{P}&0\end{pmatrix}\psi_d(x)-m_0c\psi_d(x)=0.
\end{align}

For the limit that $ m=0 $, then

\begin{align}
\bm{P}_c\psi_l=0,~~~\bm{P}\psi_r=0.
\end{align}

\noindent and $ p^0=|\vec{p}| $, which implies that

\begin{align}
(\bm{\sigma}\cdot\hat{\bm{p}})\psi_r=\psi_r,~~~
(\bm{\sigma}\cdot\hat{\bm{p}})\psi_l=-\psi_l,
\end{align}

\end{frame}
\end{document}